\chapter*{Part Synthesis: Evidence, Reliability, and Systems}
\addcontentsline{toc}{chapter}{Part Synthesis: Evidence, Reliability, and Systems}

\begin{whychaptermatters}
This part turned mature ML practice into explicit undergraduate workflow. Students now have one chain from experimental claim to reliability review to launch-or-no-launch service judgment, which is exactly where the book's problem-driven identity becomes most distinctive.
\end{whychaptermatters}

\begin{chaptersummary}
\item \textbf{Question this part taught.} What does the result justify, where can it still fail badly enough to matter, and what must be true before a model becomes a real system?
\item \textbf{Artifact chain this part added.} Experiment claim sheet $\rightarrow$ reliability review card $\rightarrow$ system readiness brief.
\item \textbf{What a student can now do.} Turn an offline result into a bounded scientific claim, test that claim against reliability hazards, and make a technical launch or no-launch recommendation for the whole service.
\end{chaptersummary}
